% Options for packages loaded elsewhere
\PassOptionsToPackage{unicode}{hyperref}
\PassOptionsToPackage{hyphens}{url}
%
\documentclass[
  11pt,
]{article}
\usepackage{amsmath,amssymb}
\usepackage{iftex}
\ifPDFTeX
  \usepackage[T1]{fontenc}
  \usepackage[utf8]{inputenc}
  \usepackage{textcomp} % provide euro and other symbols
\else % if luatex or xetex
  \usepackage{unicode-math} % this also loads fontspec
  \defaultfontfeatures{Scale=MatchLowercase}
  \defaultfontfeatures[\rmfamily]{Ligatures=TeX,Scale=1}
\fi
\usepackage{lmodern}
\ifPDFTeX\else
  % xetex/luatex font selection
\fi
% Use upquote if available, for straight quotes in verbatim environments
\IfFileExists{upquote.sty}{\usepackage{upquote}}{}
\IfFileExists{microtype.sty}{% use microtype if available
  \usepackage[]{microtype}
  \UseMicrotypeSet[protrusion]{basicmath} % disable protrusion for tt fonts
}{}
\makeatletter
\@ifundefined{KOMAClassName}{% if non-KOMA class
  \IfFileExists{parskip.sty}{%
    \usepackage{parskip}
  }{% else
    \setlength{\parindent}{0pt}
    \setlength{\parskip}{6pt plus 2pt minus 1pt}}
}{% if KOMA class
  \KOMAoptions{parskip=half}}
\makeatother
\usepackage{xcolor}
\usepackage[margin=25mm]{geometry}
\setlength{\emergencystretch}{3em} % prevent overfull lines
\providecommand{\tightlist}{%
  \setlength{\itemsep}{0pt}\setlength{\parskip}{0pt}}
\setcounter{secnumdepth}{-\maxdimen} % remove section numbering
\providecommand{\gt}{>}
\providecommand{\lt}{<}
\ifLuaTeX
  \usepackage{selnolig}  % disable illegal ligatures
\fi
\IfFileExists{bookmark.sty}{\usepackage{bookmark}}{\usepackage{hyperref}}
\IfFileExists{xurl.sty}{\usepackage{xurl}}{} % add URL line breaks if available
\urlstyle{same}
\hypersetup{
  hidelinks,
  pdfcreator={LaTeX via pandoc}}

\author{}
\date{}

\begin{document}

Exponential and logarithmic differentiation

MIT 18.01, Fall 2006, Lecture 6

This lesson develops one useful connection: taking a logarithm turns a
changing power into a product, and its derivative measures relative
change. You will derive the exponential and logarithm derivative rules,
use them when both a base and an exponent vary, and evaluate a limit by
recognising a derivative.

Reading route: read Parts 1--5 in order and try A1--A5 where they occur.
Keep any partial working; the \protect\hyperlink{hints}{hints} give a
next step and the \protect\hyperlink{solutions}{solutions} give complete
reasoning in a separate group. Use \protect\hyperlink{task-a6}{A6} on a
later return. The source and its corrections are recorded at the end.

Starting tools: the derivative is a tangent slope or instantaneous rate,
defined by a difference quotient; the product rule is \((uv)'=u'v+uv'\);
and the chain rule is \((F(u(x)))'=F'(u(x))u'(x)\). We use ordinary
algebra, limits, continuous functions, exponent laws, and the meaning of
an inverse function. The needed connections to logarithms are developed
below. All quantities in powers and logarithms are real, and logarithm
inputs must be positive.

Part 1: why an exponential's derivative is proportional to itself

Start with a fixed base \(a\gt 1\), as in the lecture. In \(a^x\), \(a\)
stays fixed while \(x\) changes. The exponent laws give

\[
a^0=1,\qquad a^1=a,\qquad a^2=a\cdot a,
\]

\[
a^{x_1+x_2}=a^{x_1}a^{x_2},\qquad (a^{x_1})^{x_2}=a^{x_1x_2}.
\]

For integers \(p\) and positive integers \(q\),
\(a^{p/q}=\sqrt[q]{a^p}\). Extending from rational exponents to real
ones by continuity gives the real exponential family. We use its usual
differentiability; constructing that family rigorously is not this
lesson's task.

Write \(h\) for a small nonzero change in \(x\). The lecture writes
\(\Delta x\) for the same change. At a fixed \(x\), the difference
quotient is

\[
\frac{a^{x+h}-a^x}{h}
=a^x\frac{a^h-1}{h}.
\]

The exponent law factors \(a^{x+h}\) as \(a^xa^h\). While \(h\) tends to
zero, \(a^x\) is constant, so it can be taken outside the limit. Define

\[
M(a)=\lim_{h\to0}\frac{a^h-1}{h}.
\]

\(M\) is a function of the base: it assigns each base its slope at
\(x=0\). It is a constant with respect to \(x\) when that base is fixed.
The definition now gives

\[
\frac{d}{dx}a^x=M(a)a^x.
\]

To read this formula, multiply the current height \(a^x\) by the base's
fixed factor \(M(a)\) to get the slope there. To construct it from the
graph, first take the tangent slope at \((0,1)\); that number is
\(M(a)\), because \(a^0=1\). Then the slope at any other \(x\) is that
number times the height. This is the relationship shown by the lecture's
first figure.

The special base \(e\)

The natural exponential uses the unique base \(e\) for which \(M(e)=1\),
equivalently

\[
\lim_{h\to0}\frac{e^h-1}{h}=1.
\]

We take the existence of this special base as a foundational property of
the real exponential family. The lecture's following geometric
comparisons motivate where it lies; pictures alone do not prove
existence or uniqueness.

\begin{itemize}
\tightlist
\item
  For \(y=2^x\), the secant through \((0,1)\) and \((1,2)\) has slope
  \((2-1)/(1-0)=1\). The upward-curving graph has a smaller tangent
  slope at the left endpoint, so \(M(2)\lt 1\).
\item
  For \(y=4^x\), the secant through \((-1/2,1/2)\) and \((0,1)\) has
  slope \((1-1/2)/(0-(-1/2))=1\). At the right endpoint the tangent is
  steeper, so \(M(4)\gt 1\).
\end{itemize}

These are the comparisons in Figures 2 and 3: a secant measures an
average slope between its two points, while a tangent measures the slope
at one point. For these strictly upward-curving exponentials, tangent
slopes increase along the curve, placing the secant slope between the
endpoint tangent slopes. The comparisons locate the slope-one base
between 2 and 4 once continuity and increase of \(M(a)\) are
established. We will identify \(M(a)\) explicitly in Part 2, which also
verifies those properties. The printed source labels the second secant's
right endpoint \((1,0)\); the correct coordinates are \((0,1)\).

Substituting \(M(e)=1\) into the derivative formula yields

\[
\frac{d}{dx}e^x=e^x.
\]

The chain rule extends this to a differentiable input \(u(x)\):

\[
\frac{d}{dx}e^{u(x)}=e^{u(x)}u'(x).
\]

For example, to differentiate \(e^{3x}\), the outer function is \(e^u\)
and the inner input is \(u=3x\). The outer derivative is evaluated at
that same input, giving \(e^{3x}\), and the inner derivative is 3. Thus
\((e^{3x})'=3e^{3x}\). Changing how quickly the exponent changes changes
how quickly the whole expression changes.

Task A1 (generated)

Let \(y=e^{2x-1}\). Find \(dy/dx\) and the tangent slope at \(x=1/2\).
Explain where the factor from the inner expression enters. A complete
response connects the chain rule to the factor and evaluates the
derivative at the stated point.

\protect\hyperlink{hint-a1}{Hint A1} ·
\protect\hyperlink{solution-a1}{Solution A1}

Part 2: logarithms reverse exponentials

The natural logarithm answers an inverse question: which exponent of
\(e\) gives a specified positive number? Thus

\[
y=e^x\quad\Longleftrightarrow\quad\ln y=x,
\]

and, with different variable names,

\[
w=\ln x\quad\Longleftrightarrow\quad e^w=x\qquad(x\gt 0).
\]

For instance, \(\ln(e^3)=3\) and \(e^{\ln 5}=5\). The notation \(\ln x\)
is a function value, not a product of symbols. Since \(e^0=1\) and the
exponential is increasing, \(\ln1=0\), \(\ln x\lt 0\) for
\(0\lt x\lt 1\), and \(\ln x\gt 0\) for \(x\gt 1\). The logarithm has no
real value at zero or at a negative input.

For positive \(u,v\), the exponent law gives the logarithm law

\[
\ln(uv)=\ln u+\ln v.
\]

Indeed, put \(u=e^r\) and \(v=e^s\). Then \(uv=e^{r+s}\), so its
logarithm is \(r+s\). Similarly, for \(u\gt 0\) and real \(r\),
\(\ln(u^r)=r\ln u\). Moving an exponent outside the logarithm will be
our main simplifying step.

To obtain the logarithm's derivative, set \(w(x)=\ln x\), so
\(e^{w(x)}=x\). The positive exponential derivative means the inverse
has a finite local derivative here. Differentiate this identity with
respect to \(x\) using the chain rule:

\[
e^{w(x)}w'(x)=1.
\]

Since \(e^{w(x)}=x\gt 0\), division is permitted, giving

\[
\frac{d}{dx}\ln x=\frac1x\qquad(x\gt 0).
\]

The derivative is positive, confirming that \(\ln x\) is increasing;
differentiability also gives continuity on the positive inputs. For a
positive differentiable function \(u(x)\), the same chain rule gives

\[
\frac{d}{dx}\ln(u(x))=\frac{u'(x)}{u(x)}.
\]

The numerator is the derivative of the entire inner input; the
denominator is that input itself. For example, \(\ln(1+x^2)\) is defined
for every real \(x\), and its derivative is \(2x/(1+x^2)\). A check at
\(x=0\) gives zero slope, consistent with the expression being unchanged
when \(x\) is replaced by \(-x\) and with the displayed derivative.

Returning to a general constant base

For fixed \(a\gt 0\), write \(a=e^{\ln a}\). The exponent laws then give

\[
a^x=(e^{\ln a})^x=e^{x\ln a}.
\]

This rewrites the same quantity in base \(e\); it does not change its
value or interchange the base and exponent. Since \(\ln a\) is constant
with respect to \(x\), the chain rule gives

\[
\frac{d}{dx}a^x=(\ln a)e^{x\ln a}=(\ln a)a^x.
\]

Comparison with Part 1 identifies \(M(a)=\ln a\). In particular \(M\) is
continuous and strictly increasing for \(a\gt 0\), and \(M(a)=1\) has
the unique solution \(a=e\). The previous geometric bracket
\(2\lt e\lt 4\) is consistent with this identification. This is a check
of the slope-based characterisation, not a new construction of \(e\)
from the pictures.

Even base 10 brings in the natural logarithm: \((10^x)'=(\ln10)10^x\).
The same argument extends the lecture's initial \(a\gt 1\) convention to
\(0\lt a\lt 1\): there \(\ln a\lt 0\), so \(a^x\) decreases. At \(a=1\)
the function is constant and the derivative is zero. No corresponding
general real-variable rule with a negative base is being claimed.

Task A2 (generated)

Find the derivatives of \(f(x)=(1/2)^{3x}\) and \(g(x)=\ln(5-2x)\).
State the real domain of each function and explain why each derivative
is negative throughout its domain. A complete response includes the
inner derivative, the domain inequality for the logarithm, and a sign
justification.

\protect\hyperlink{hint-a2}{Hint A2} ·
\protect\hyperlink{solution-a2}{Solution A2}

Part 3: logarithmic differentiation when a power changes

Let \(f(x)\gt 0\) be differentiable. The previous chain-rule result says

\[
(\ln f(x))'=\frac{f'(x)}{f(x)},\qquad
f'(x)=f(x)(\ln f(x))'.
\]

A prime means differentiation with respect to the current variable, here
\(x\). In particular \((\ln f)'\) is the derivative of the composite
function \(x\mapsto\ln(f(x))\). It is not \(\ln(f')\). The first formula
divides absolute change rate by the current value; the second recovers
the absolute rate by multiplying back by that value.

This gives a procedure: take the logarithm of a positive expression,
simplify using logarithm laws, differentiate the simplified expression,
and multiply by the original function. Positivity is a condition for
taking the real logarithm in this procedure, not a claim that every
differentiable function is positive.

Applied to \(f(x)=a^x\), it provides the lecture's second method:

\[
\ln f=x\ln a,\qquad (\ln f)'=\ln a,
\]

\[
f'=f\ln a=(\ln a)a^x.
\]

Here \(\ln a\) is fixed. For this short expression, rewriting in base
\(e\) is equally effective.

Now consider the lecture's worked example \(f(x)=x^x\) on \(x\gt 0\).
Both the base and exponent change. The ordinary power rule
\((x^n)'=nx^{n-1}\) treats \(n\) as constant, while the exponential rule
\((a^x)'=(\ln a)a^x\) treats \(a\) as constant. Neither rule alone
applies to the two changing occurrences of \(x\).

Taking the logarithm exposes those two changes as a product:

\[
\ln f=x\ln x.
\]

The product rule now gives

\[
(\ln f)'=1\cdot\ln x+x\cdot\frac1x=\ln x+1.
\]

Multiplying by the original function recovers

\[
f'(x)=x^x(\ln x+1)\qquad(x\gt 0).
\]

For example, at \(x=2\) the slope is \(4(\ln2+1)\), not 4 from treating
the exponent as fixed at 2. Alternatively, \(x^x=e^{x\ln x}\), and the
chain rule gives precisely the same result. These two methods use the
same simplification of the changing power.

More generally, if \(u(x)\gt 0\) and both \(u,v\) are differentiable,
construct the real power \(y=u(x)^{v(x)}\) as \(e^{v(x)\ln u(x)}\). Then

\[
\ln y=v\ln u,\qquad
\frac{y'}{y}=v'\ln u+v\frac{u'}u.
\]

The first term records the changing exponent; the second records the
changing base. Setting \(v\) constant removes the first term and
recovers the ordinary power rule with the inner derivative. Setting
\(u\) constant removes the second and recovers the constant-base
exponential rule. This connects the methods rather than creating an
unrelated formula to memorise.

Task A3 (generated)

For \(y=(1+x)^{x^2}\), \(x\gt -1\), find \(y'\) and the tangent slope at
\(x=0\). Explain why treating the exponent \(x^2\) as constant
throughout differentiation would miss a term. A complete response
differentiates both changing parts after a justified rewrite and
preserves the given domain.

\protect\hyperlink{hint-a3}{Hint A3} ·
\protect\hyperlink{solution-a3}{Solution A3}

Part 4: a limit hidden inside a logarithm

The lecture asks for the limit of

\[
b_k=\left(1+\frac1k\right)^k\qquad(k=1,2,3,\ldots).
\]

The subscript labels the term of a sequence; \(k\to\infty\) means
considering larger and larger positive integers. Although the base tends
to 1, the exponent grows. Replacing the base by its limit too early
loses the effect of many repeated factors. Nor can the competing factors
\(k\) and \(\ln(1+1/k)\) be evaluated separately as an ordinary product
of finite limits.

Each \(b_k\) is positive, so define \(a_k=\ln b_k\). These subscripted
letters label sequences; they are not the fixed base \(a\) from Part 1.
The logarithm law gives

\[
a_k=k\ln\left(1+\frac1k\right)
=\frac{\ln(1+1/k)}{1/k}.
\]

The change of variable \(h=1/k\) makes \(h\to0\) from the positive side.
Using \(\ln1=0\) gives the exact equality

\[
a_k=\frac{\ln(1+h)-\ln1}{h}.
\]

Compare this with the derivative definition \([F(x_0+h)-F(x_0)]/h\):
here \(F=\ln\) and \(x_0=1\). Thus Part 2 supplies the limit

\[
a_k\longrightarrow (\ln x)'\big|_{x=1}=\frac11=1.
\]

The vertical bar means evaluate the derivative at the stated input. We
have found the logarithm's limit. To recover the original sequence, use
the exact inverse relation \(b_k=e^{a_k}\) and continuity of the
exponential:

\[
b_k=e^{a_k}\longrightarrow e^1=e.
\]

Consequently,

\[
\lim_{k\to\infty}\left(1+\frac1k\right)^k=e.
\]

We did not assume that \(b_k\) already converges in order to take the
logarithm of its proposed limit. Instead, we found a convergent sequence
\(a_k\) and then applied a continuous function to it.

The lecture's first closing remark uses this formula to approximate
\(e\). With \(k=10\), it gives \((1.1)^{10}\approx2.59374\), whereas
\(e\approx2.71828\): the estimate is about 4.58\% low. It is useful as a
rough approximation, and the limit guarantees convergence as \(k\)
increases; \(k=10\) does not give several correct decimal places.

The same route works with a small negative increment as long as the
logarithm's input stays positive: the derivative of \(\ln\) at 1 has the
same limit from either side. More broadly, if logarithms tend to a
finite \(L\), exponentiating yields a limit \(e^L\). If logarithms
instead tend to \(+\infty\), the original positive quantities grow
without bound because the exponential is increasing and unbounded. If
they tend to \(-\infty\), the original quantities tend to zero.

Task A4 (generated)

For positive integers \(k\gt 2\), let \(b_k=(1-2/k)^{3k}\). Find the
limit of \(\ln(b_k)\), then the limit of \(b_k\). Show how \(h=-2/k\)
converts \(\ln(b_k)\) into a multiple of the derivative quotient for
\(\ln\) at 1, and state the side from which \(h\) approaches zero. A
complete response distinguishes the logarithm's limit from the original
sequence's limit and justifies returning by the exponential.

\protect\hyperlink{hint-a4}{Hint A4} ·
\protect\hyperlink{solution-a4}{Solution A4}

Part 5: relative change explains the usefulness of logarithms

A fall of 50 points has different significance from a starting level of
300 than from 10000. For positive initial and final levels \(F_0,F_1\),
the exact finite fractional change is

\[
\frac{F_1-F_0}{F_0},
\]

and multiplying by 100 expresses it as a percentage. The denominator is
the starting level, not the final level. This is the distinction behind
the lecture's final market example; no actual market observation is
needed.

For a positive differentiable level \(F(t)\), its instantaneous relative
rate is

\[
\frac{F'(t)}{F(t)}=\frac{d}{dt}\ln(F(t)).
\]

The prime now denotes differentiation with respect to time \(t\). If
\(F\) is measured in points and time in days, \(F'\) has units points
per day and \(F'/F\) has units per day. Multiplying by 100 gives an
instantaneous percentage rate per day. For a physical measurement, read
the logarithm as \(\ln(F/F_{\rm ref})\) with a fixed positive reference
in the same units; its derivative is still \(F'/F\), because the
reference is constant. This avoids taking a logarithm of a dimensional
quantity.

Why is that an instantaneous percentage rate? The derivative definition
gives, for small time increments \(\Delta t\),

\[
F(t+\Delta t)-F(t)\approx F'(t)\Delta t,
\]

and division by \(F(t)\gt 0\) yields

\[
\frac{F(t+\Delta t)-F(t)}{F(t)}
\approx\frac{F'(t)}{F(t)}\Delta t.
\]

This is a local approximation, not an exact finite-change rule. The
exact finite logarithmic change is \(\ln(F_1/F_0)\), which also
generally differs from \((F_1-F_0)/F_0\); the two are close for a small
fractional change because \((\ln x)'\) at 1 is 1.

For a worked model, let \(F(t)=80e^{0.03t}\), where \(t\) is the
numerical time in days and \(F\) is in points. Its absolute rate is
\(F'=2.4e^{0.03t}\) points per day. Dividing by its current level gives
\(F'/F=0.03\) per day, or 3\% per day instantaneously. After one day,
however, the exact percentage increase is \(100(e^{0.03}-1)\), about
3.045\%. The level on which growth acts changes during the day. Calling
both numbers exactly 3\% would confuse a local rate with a finite
change.

Task A5 (generated mathematical models, not market data)

Two positive index levels are modelled by \(F(t)=300e^{-0.02t}\) and
\(G(t)=10000e^{-0.02t}\), where \(t\) is the numerical time in days. At
\(t=0\), find each absolute rate in points per day and each
instantaneous percentage rate per day. Separately compare the exact
finite percentage changes caused by a fall of 50 points from 300 and
from 10000. For the model \(F\), express its exact percentage change
from day 0 to day 1 and explain why it differs from the instantaneous
percentage rate multiplied by one day. A complete response states which
quantities agree, carries signs and units, and separates exact finite
change from local approximation; a symbolic expression in \(e\) is
sufficient.

\protect\hyperlink{hint-a5}{Hint A5} ·
\protect\hyperlink{solution-a5}{Solution A5}

Later return: Task A6 (generated)

Return after a gap, perhaps the next day; adjust the gap to your needs.
The first request below retrieves the main relationships. The later
requests test using them to choose methods and handle a changed limit.
Needing to reopen the lesson identifies a useful place to revisit; it is
not evidence about your overall ability.

First, with the lesson closed, write what \(M(a)\) means geometrically
and give the relationship between \((\ln f)'\) and \(f'\) for positive
differentiable \(f\). Then reopen the lesson if needed and differentiate
\(p(x)=x^3\), \(q(x)=3^x\) and \(r(x)=x^x\) on \(x\gt 0\), giving a
reason for the method chosen in each case. Finally decide whether
\(c_k=(1+1/k)^{k^2}\), for positive integers \(k\), has a finite limit,
and justify the conclusion through its logarithm rather than the
misleading substitution 1 raised to an increasing power. A complete
response separates recall from transfer, identifies what is constant in
each derivative, and controls the growing factor in \(\ln(c_k)\).

\protect\hyperlink{hint-a6}{Hint A6} ·
\protect\hyperlink{solution-a6}{Solution A6}

Hints

These hints are grouped separately from the complete solutions. Stop at
whichever hint lets you resume your working.

Hint A1

Put \(u=2x-1\). Write the outer derivative as \(e^u\), find \(u'\), and
only then substitute \(x=1/2\) into their product.

\protect\hyperlink{task-a1}{Return to A1}

Hint A2

Write the first function as \(e^{3x\ln(1/2)}\). For the second, solve
\(5-2x\gt 0\) before differentiating. Its derivative has the inner
derivative in the numerator and the positive input in the denominator.

\protect\hyperlink{task-a2}{Return to A2}

Hint A3

The given interval makes \(1+x\) positive. Begin with
\(\ln y=x^2\ln(1+x)\). In the product rule, keep one contribution from
\((x^2)'\) and another from \((\ln(1+x))'\), then multiply by \(y\).

\protect\hyperlink{task-a3}{Return to A3}

Hint A4

After taking the logarithm, the expression is \(3k\ln(1-2/k)\). Since
\(h=-2/k\), solve for \(k\) as \(-2/h\). This exposes a constant
multiplying \([\ln(1+h)-\ln1]/h\).

\protect\hyperlink{task-a4}{Return to A4}

Hint A5

For either exponential model, first compute its derivative and then
divide by the same model's current value. For the separate 50-point
falls, use \(\Delta F/F_0\) with \(\Delta F=-50\). For the one-day model
change, find \(F(1)/F(0)\) before subtracting 1.

\protect\hyperlink{task-a5}{Return to A5}

Hint A6

The point used to define \(M(a)\) is \((0,1)\). For the three
derivatives, identify whether the base, exponent, or both change. For
the last part write \(\ln c_k=k[k\ln(1+1/k)]\); Part 4 already
establishes the limit of the bracketed factor. A factor tending to 1 is
eventually greater than \(1/2\).

\protect\hyperlink{task-a6}{Return to A6}

Complete solutions

Compare the first step where your working differs. A missing chain-rule
factor, an invalid logarithm input, and confusing a logarithm's limit
with the original limit are different problems. The explanations below
describe possible errors, not errors you have been observed making.

Solution A1

The outer function \(e^u\) has derivative \(e^u\), evaluated at
\(u=2x-1\). The inner derivative is 2, so

\[
\frac{dy}{dx}=2e^{2x-1}.
\]

At \(x=1/2\) the exponent is zero and \(e^0=1\), giving slope 2. The
extra factor records that a change in \(x\) changes the exponent at
twice that rate. Substituting the point before differentiating would
erase the varying function and leave only a number; differentiate first.

\protect\hyperlink{task-a1}{Return to A1} ·
\protect\hyperlink{hint-a1}{Hint A1}

Solution A2

The positive fixed base \(1/2\) permits every real \(x\). Rewriting
\(f=e^{3x\ln(1/2)}\) and applying the chain rule gives

\[
f'(x)=3\ln(1/2)(1/2)^{3x}.
\]

Here \(\ln(1/2)\lt 0\), while 3 and the exponential are positive. Hence
\(f'\lt 0\) for all real \(x\).

The second function requires \(5-2x\gt 0\), or \(x\lt 5/2\). On that
interval,

\[
g'(x)=\frac{-2}{5-2x}\lt 0.
\]

Its numerator is negative and its denominator positive. Merely writing
\(1/(5-2x)\) would omit the changing input's derivative and predict the
wrong sign. The logarithm itself is undefined at \(x=5/2\) and beyond,
so the derivative formula supplies no real-function derivative there.

\protect\hyperlink{task-a2}{Return to A2} ·
\protect\hyperlink{hint-a2}{Hint A2}

Solution A3

For \(x\gt -1\), the base is positive and so is \(y=e^{x^2\ln(1+x)}\).
Taking its logarithm is valid:

\[
\ln y=x^2\ln(1+x).
\]

Use the product and chain rules:

\[
\frac{y'}y=2x\ln(1+x)+\frac{x^2}{1+x}.
\]

Thus

\[
y'=(1+x)^{x^2}\left(2x\ln(1+x)+\frac{x^2}{1+x}\right),\qquad x\gt -1.
\]

At \(x=0\), \(y=1\) and both terms in parentheses vanish, giving slope
0. The term \(2x\ln(1+x)\) comes from the changing exponent. Treating
\(x^2\) as constant would lose that term. Evaluating only at zero would
hide that particular mistake because the missing term happens to vanish
there; the full derivative must still include it.

\protect\hyperlink{task-a3}{Return to A3} ·
\protect\hyperlink{hint-a3}{Hint A3}

Solution A4

For \(k\gt 2\), \(1-2/k\gt 0\), so the real logarithm is allowed. Taking
it gives

\[
\ln b_k=3k\ln(1-2/k).
\]

With \(h=-2/k\), \(h\to0\) from below and \(k=-2/h\). Consequently

\[
\ln b_k=-6\frac{\ln(1+h)-\ln1}{h}\longrightarrow-6.
\]

The quotient tends to \((\ln x)'\) at 1, which is 1, just as in Part 4;
approaching from below is permitted. Finally \(b_k=e^{\ln b_k}\), so
continuity yields

\[
b_k\longrightarrow e^{-6}.
\]

A result of \(-6\) for the original limit would miss this last inverse
operation and would also conflict with positivity of every \(b_k\).

\protect\hyperlink{task-a4}{Return to A4} ·
\protect\hyperlink{hint-a4}{Hint A4}

Solution A5

By the chain rule,

\[
F'(t)=-0.02F(t),\qquad G'(t)=-0.02G(t).
\]

At zero, the absolute rates are \(-0.02(300)=-6\) points per day and
\(-0.02(10000)=-200\) points per day. They differ because the starting
levels differ. Dividing by the respective levels gives the same relative
rate \(-0.02\) per day, or \(-2\text{ percent}\) per day
instantaneously.

For the separate finite falls, the exact percentages are

\[
\frac{-50}{300}\times100\text{ percent}=-\frac{50}{3}\text{ percent}\approx-16.67\text{ percent},\qquad
\frac{-50}{10000}\times100\text{ percent}=-0.5\text{ percent}.
\]

These are changes, not rates, because no elapsed time is specified for
those falls. They are comparisons separate from the two exponential
models.

For \(F\) over one day, \(F(1)/F(0)=e^{-0.02}\), so its exact percentage
change is

\[
100(e^{-0.02}-1)\text{ percent}\approx-1.9801\text{ percent}.
\]

The initial relative rate times one day gives \(-2\text{ percent}\), a
local approximation. As the level falls, its absolute loss rate also
becomes smaller in magnitude. The exact one-day fall is therefore
slightly smaller than 2\% of its starting level. Equal instantaneous
percentage rates do not make finite percentages exactly equal to the
rate times elapsed time.

\protect\hyperlink{task-a5}{Return to A5} ·
\protect\hyperlink{hint-a5}{Hint A5}

Solution A6

The recall targets are the tangent slope of \(y=a^x\) at \((0,1)\),
called \(M(a)\), and \((\ln f)'=f'/f\) for positive differentiable
\(f\). Equivalently \(f'=f(\ln f)'\). Also \(M(a)=\ln a\).

For \(p=x^3\), the exponent is the fixed number 3, so the power rule
gives \(p'=3x^2\). For \(q=3^x\), the base is fixed, so
\(q'=(\ln3)3^x\). For \(r=x^x\), both parts change; logarithmic
differentiation or rewriting as \(e^{x\ln x}\) gives
\(r'=x^x(\ln x+1)\). Alternative correct methods are acceptable if they
account for every changing part and preserve the domain \(x\gt 0\).

For the sequence, take a logarithm:

\[
\ln c_k=k^2\ln(1+1/k)
=k\left[k\ln(1+1/k)\right].
\]

The bracketed factor tends to 1 by Part 4. To make the conclusion
precise, it is greater than \(1/2\) for all sufficiently large \(k\), so
\(\ln c_k\gt k/2\) eventually. Thus \(\ln c_k\to+\infty\), and
\(c_k\to+\infty\) by the exponential's increase without bound. There is
no finite limit. The additional factor \(k\) in the exponent changes the
conclusion from the lecture's \(e\) limit; the base approaching 1 does
not determine the answer by itself.

\protect\hyperlink{task-a6}{Return to A6} ·
\protect\hyperlink{hint-a6}{Hint A6}

Source and scope

This lesson covers the complete supplied
\href{https://ocw.mit.edu/courses/18-01-single-variable-calculus-fall-2006/f9af0e98490296c99d330faf47389507_lec6.pdf}{MIT
OpenCourseWare 18.01 Fall 2006 Lecture 6 PDF}, comprising a cover and
printed pages 1--7. The printed title is ``Exponential and Log,
Logarithmic Differentiation, Hyperbolic Functions.'' Hyperbolic
functions appear in that title but are not developed anywhere in this
PDF; no missing continuation is assumed here.

The source's worked \(x^x\) derivative, exponential limit, three
geometric figures and both final remarks are developed above. Tasks
A1--A6 are generated applications, not quoted MIT assessment questions.
Two source typos are corrected: the point labelled \((1,0)\) in the
discussion of \(4^x\) and Figure 3 is \((0,1)\), and the inverse box on
printed page 4 must read \(w=\ln x\Rightarrow e^w=x\). The approximation
at \(k=10\) is quantified in Part 4, and finite percentage change is
distinguished from instantaneous relative rate in Part 5.

\end{document}
